\chapter{Profilo degli score attorno ai cluster concordati}
\label{app:profilo}

Questa appendice mostra, come sola visualizzazione, come cambia lo score di ciascun metodo dall'interno dei cluster concordati verso l'esterno (Figura~\ref{fig:app-profilo}). In ogni istante degli anni di test ogni cella ha una posizione, contata in celle della griglia dal bordo del cluster concordato più vicino, con le celle in diagonale a distanza uno (una cella misura circa 55~km in latitudine e meno in longitudine): negativa dentro il cluster, positiva fuori; l'ultima posizione raccoglie tutte le celle a sei o più celle dal bordo. Per ogni posizione la figura riporta il percentile mediano di ciascuno score (Sezione~\ref{sec:exp-misure}) fra i suoi valori degli anni di test, separatamente per i cluster concordati brevi e lunghi; le celle più vicine a un cluster concordato piccolo non sono riportate. Dentro i valori sono alti per costruzione, perché almeno cinque metodi su nove superano lì la propria soglia e il bordo è definito dagli stessi score; fuori il calo riflette anche la continuità nello spazio dei campi. A una cella dal bordo i percentili mediani sono fra 82 e 89 nei cluster brevi e fra 78 e 88 nei lunghi, a sei o più celle fra 36 e 41. I profili dei nove score sono quasi sovrapposti: nessuno score si comporta in modo anomalo attorno ai cluster concordati. È un controllo di coerenza, non un confronto fra i metodi.
% Dati: figures/dati/profilo_r6.csv, da scripts/check_era5_event_profile.py --cube outputs/concordate_5pc_regola_tesi
% --inputs-run outputs/stgan_optuna27_scored --limit 6 (physiq-pv a7dd5d69, anni di test 2005-2025; conteggi dei cluster
% uguali al cubo). Posizione +1: brevi 82 (variazione) - 89 (G_n), lunghi 78 (variazione) - 88 (G_n, PCA, COPOD, iForest);
% posizione >= 6: brevi 36 (G_n) - 41 (variazione), lunghi 38 (G_n) - 41 (CBLOF, normale, variazione). Figura:
% figures/genera_figura_profilo_concordate.py. Inventario delle fonti: classe C, visualizzazione senza fonte, per scelta
% dell'autore (16 novembre 2026).

\begin{figure}[!htb]
    \centering
    \includegraphics[width=\textwidth]{era5_concordate_profilo}
    \caption{Percentile mediano di ogni score secondo la distanza in celle dal bordo del cluster concordato più vicino, nei cluster brevi (a sinistra) e lunghi (a destra), sugli anni di test; i punti a destra raccolgono le celle a sei o più celle. Dentro il cluster almeno cinque metodi su nove superano la propria soglia per costruzione. Le tre curve e la fascia dei sei modelli shallow sono quasi sovrapposte.}
    \label{fig:app-profilo}
\end{figure}
